% GENERATED by make_ablation_tables.py -- do not edit.
% Regenerate with: python make_ablation_tables.py
\begin{table}[htbp]
\centering
\small
\setlength{\tabcolsep}{5pt}
\caption{The large-season classifier at both aggregation levels, five
evaluation seasons 2021--2025. The figure in parentheses is the number of
states, respectively municipalities, that enter the level; the number
evaluated per season is smaller, since a place without the required
history is excluded rather than imputed (Section~\ref{sec:targets}).
The target is whether a unit exceeds the 75th
percentile of its own history; the loss is the Brier score. Positive values
mean the later rung has the lower Brier score, and ``median'' is over the
five seasons. ``Seasons CI excl.\ 0'' counts the seasons with a
same-direction improvement whose 95\% unit-bootstrap interval excludes
zero, which is what the decision rule of Section~\ref{sec:ablation}
requires in at least 3 of 5. The municipal rows use a within-municipality
threshold, without which the label is a population indicator; there the
population-only reference reaches a median AUC of $0.47$ against $0.58$--$0.65$ for the fitted rungs.}
\label{tab:exceed}
\begin{tabular}{@{}llrccr@{}}
\toprule
Level & Contrast & \makecell{$\Delta$Brier \\ median} & Seasons $\Delta>0$ & \makecell{Seasons \\ CI excl.\ 0} & Verdict\\
\midrule
\makecell{State \\ ($n=26$)} & \makecell{M1$-$M0 \\ (lagged metrics)} & $+0.126$ & 4/5 & 2/5 & not supported\\
 & \makecell{M2$-$M1 \\ (climate)} & $+0.027$ & 3/5 & 2/5 & not supported\\
 & \makecell{M3$-$M2 \\ (scaling)} & $+0.000$ & 3/5 & 0/5 & not supported\\
\addlinespace
\makecell{Municipal \\ ($n=3{,}976$)} & \makecell{M1$-$M0 \\ (lagged metrics)} & $+0.049$ & 4/5 & 4/5 & \textbf{supported}\\
 & \makecell{M2$-$M1 \\ (climate)} & $+0.003$ & 4/5 & 2/5 & not supported\\
 & \makecell{M3$-$M2 \\ (scaling)} & $+0.002$ & 5/5 & 1/5 & not supported\\
\bottomrule
\end{tabular}
\end{table}
